% GENERATED FILE. Edit canonical lessons, then run npm run generate:books.
% canonical-source-hash: fnv1a64:511954a8cc9ee8f1
% canonical-lessons: UR-C363-andesha, UR-C363-ishara, UR-C363-alamat, UR-C363-khadsha, UR-C363-etimad

\chapter{Hopes and Fears}
\label{ch:ur-hopes-and-fears}
\hlchaptermodality{363}

\begin{chapteropening}
\textbf{By the end of this chapter:} \emph{I can say what I hope for and what I fear might happen.}

Complete the last lesson of chapter 363: I can say what I hope for and what I fear might happen.
\end{chapteropening}

\section[\texorpdfstring{\ur{اندیشہ} (andeshā)}{اندیشہ (andeshā)}]{\ur{اندیشہ} (andeshā) — an apprehension; a misgiving}

\label{lesson:UR-C363-andesha}



\begin{quote}
Before the new one: say the Urdu for conjecture, then the Urdu for to work out a riddle.
\end{quote}

\subsection*{You'll want to know: \ur{اندیشہ}}
\textbf{\ur{اندیشہ}} — \emph{andeshā} — \textquotedblleft{}an apprehension; a misgiving\textquotedblright{}.

\begin{quote}
\ur{اگر} \ur{بارش} \ur{ہوئی} \ur{تو} \ur{سیلاب} \ur{کا} \ur{اندیشہ} \ur{ہے۔}
\end{quote}

\begin{quote}
\emph{agar bārish huī to sailāb kā andeshā hai.}
\end{quote}

\begin{quote}
If it rains, there is a fear of floods.
\end{quote}

Masculine.

\begin{cousinweb}
\textbf{\ur{اندیشہ}} (\emph{andeshā}) is usually traced to Persian \emph{andēsha}, \textquotedblleft{}thought, anxiety\textquotedblright{}: a thought that worries.
\end{cousinweb}

\begin{grammarlens}[title={What you've built}]
One more word for saying what happens if.
\end{grammarlens}

\subsection*{Your turn}
\begin{itemize}
\raggedright
  \item \emph{Say it:} \emph{andeshā}
  \item \emph{Say it:} \emph{andeshā}, once more
  \item \emph{Say it:} \emph{agar bārish huī to sailāb kā andeshā hai.} — If it rains, there is a fear of floods.
  \item \emph{Recall:} say the Urdu for conjecture, then the Urdu for to work out a riddle, then say \emph{andeshā} again
\end{itemize}

\begin{tcolorbox}[breakable,colback=teal!4,colframe=teal!35!black,title={Before you move on}]
What does \ur{اندیشہ} mean? (\textquotedblleft{}An apprehension; a misgiving\textquotedblright{}.) Say it once more.
\end{tcolorbox}
\section[\texorpdfstring{\ur{اشارہ} (ishārā)}{اشارہ (ishārā)}]{\ur{اشارہ} (ishārā) — a hint; a signal}

\label{lesson:UR-C363-ishara}



\begin{quote}
Before the new one: say the Urdu for to work out a riddle, then the Urdu for an apprehension.
\end{quote}

\subsection*{You'll want to know: \ur{اشارہ}}
\textbf{\ur{اشارہ}} — \emph{ishārā} — \textquotedblleft{}a hint; a signal\textquotedblright{}.

\begin{quote}
\ur{اگر} \ur{مدد} \ur{چاہیے} \ur{تو} \ur{مجھے} \ur{اشارہ} \ur{دیجیے۔}
\end{quote}

\begin{quote}
\emph{agar madad chāhiye to mujhe ishārā dījiye.}
\end{quote}

\begin{quote}
If you need help, give me a sign.
\end{quote}

Masculine.

\begin{cousinweb}
\textbf{\ur{اشارہ}} (\emph{ishārā}) is usually traced, through Persian, to Arabic \emph{ishāra}, \textquotedblleft{}a sign, a gesture\textquotedblright{}.
\end{cousinweb}

\begin{grammarlens}[title={What you've built}]
One more word for saying what happens if.
\end{grammarlens}

\subsection*{Your turn}
\begin{itemize}
\raggedright
  \item \emph{Say it:} \emph{ishārā}
  \item \emph{Say it:} \emph{ishārā}, once more
  \item \emph{Say it:} \emph{agar madad chāhiye to mujhe ishārā dījiye.} — If you need help, give me a sign.
  \item \emph{Recall:} say the Urdu for to work out a riddle, then the Urdu for an apprehension, then say \emph{ishārā} again
\end{itemize}

\begin{tcolorbox}[breakable,colback=teal!4,colframe=teal!35!black,title={Before you move on}]
What does \ur{اشارہ} mean? (\textquotedblleft{}A hint; a signal\textquotedblright{}.) Say it once more.
\end{tcolorbox}
\section[\texorpdfstring{\ur{علامت} (alāmat)}{علامت (alāmat)}]{\ur{علامت} (alāmat) — a symptom; a sign (of something)}

\label{lesson:UR-C363-alamat}



\begin{quote}
Before the new one: say the Urdu for an apprehension, then the Urdu for a hint.
\end{quote}

\subsection*{You'll want to know: \ur{علامت}}
\textbf{\ur{علامت}} — \emph{alāmat} — \textquotedblleft{}a symptom; a sign (of something)\textquotedblright{}.

\begin{quote}
\ur{اگر} \ur{بخار} \ur{ہے} \ur{تو} \ur{یہ} \ur{بیماری} \ur{کی} \ur{علامت} \ur{ہے۔}
\end{quote}

\begin{quote}
\emph{agar bukhār hai to yih bīmārī kī alāmat hai.}
\end{quote}

\begin{quote}
If there is a fever, it is a sign of illness.
\end{quote}

Feminine.

\begin{cousinweb}
\textbf{\ur{علامت}} (\emph{alāmat}) is usually traced, through Persian, to Arabic \emph{ʿalāma}, \textquotedblleft{}a mark, a sign\textquotedblright{}.
\end{cousinweb}

\begin{grammarlens}[title={What you've built}]
One more word for saying what happens if.
\end{grammarlens}

\subsection*{Your turn}
\begin{itemize}
\raggedright
  \item \emph{Say it:} \emph{alāmat}
  \item \emph{Say it:} \emph{alāmat}, once more
  \item \emph{Say it:} \emph{agar bukhār hai to yih bīmārī kī alāmat hai.} — If there is a fever, it is a sign of illness.
  \item \emph{Recall:} say the Urdu for an apprehension, then the Urdu for a hint, then say \emph{alāmat} again
\end{itemize}

\begin{tcolorbox}[breakable,colback=teal!4,colframe=teal!35!black,title={Before you move on}]
What does \ur{علامت} mean? (\textquotedblleft{}A symptom; a sign (of something)\textquotedblright{}.) Say it once more.
\end{tcolorbox}
\section[\texorpdfstring{\ur{خدشہ} (khadshā)}{خدشہ (khadshā)}]{\ur{خدشہ} (khadshā) — a nagging fear (that something may happen)}

\label{lesson:UR-C363-khadsha}



\begin{quote}
Before the new one: say the Urdu for a hint, then the Urdu for a symptom.
\end{quote}

\subsection*{You'll want to know: \ur{خدشہ}}
\textbf{\ur{خدشہ}} — \emph{khadshā} — \textquotedblleft{}a nagging fear (that something may happen)\textquotedblright{}.

\begin{quote}
\ur{اگر} \ur{دیر} \ur{ہوئی} \ur{تو} \ur{گاڑی} \ur{چھوٹنے} \ur{کا} \ur{خدشہ} \ur{ہے۔}
\end{quote}

\begin{quote}
\emph{agar der huī to gāṛī chhūṭne kā khadshā hai.}
\end{quote}

\begin{quote}
If we are late, we may miss the train.
\end{quote}

Masculine.

\begin{cousinweb}
\textbf{\ur{خدشہ}} (\emph{khadshā}) is usually traced, through Persian, to Arabic \emph{khadsha}, \textquotedblleft{}a scratch\textquotedblright{}; Urdu uses it for a fear that something may go wrong, a sense our sources do not trace.
\end{cousinweb}

\begin{grammarlens}[title={What you've built}]
One more word for saying what happens if.
\end{grammarlens}

\subsection*{Your turn}
\begin{itemize}
\raggedright
  \item \emph{Say it:} \emph{khadshā}
  \item \emph{Say it:} \emph{khadshā}, once more
  \item \emph{Say it:} \emph{agar der huī to gāṛī chhūṭne kā khadshā hai.} — If we are late, we may miss the train.
  \item \emph{Recall:} say the Urdu for a hint, then the Urdu for a symptom, then say \emph{khadshā} again
\end{itemize}

\begin{tcolorbox}[breakable,colback=teal!4,colframe=teal!35!black,title={Before you move on}]
What does \ur{خدشہ} mean? (\textquotedblleft{}A nagging fear (that something may happen)\textquotedblright{}.) Say it once more.
\end{tcolorbox}
\section[\texorpdfstring{\ur{اعتماد} (etimād)}{اعتماد (etimād)}]{\ur{اعتماد} (etimād) — confidence; trust}

\label{lesson:UR-C363-etimad}



\begin{quote}
Before the new one: say the Urdu for a symptom, then the Urdu for a nagging fear.
\end{quote}

\subsection*{You'll want to know: \ur{اعتماد}}
\textbf{\ur{اعتماد}} — \emph{etimād} — \textquotedblleft{}confidence; trust\textquotedblright{}.

\begin{quote}
\ur{اگر} \ur{اعتماد} \ur{ہے} \ur{تو} \ur{آپ} \ur{ضرور} \ur{جیتیں} \ur{گے۔}
\end{quote}

\begin{quote}
\emph{agar etimād hai to āp zarūr jīteṅ ge.}
\end{quote}

\begin{quote}
If you have confidence, you will certainly win.
\end{quote}

Masculine.

\begin{cousinweb}
\textbf{\ur{اعتماد}} (\emph{etimād}) is usually traced, through Persian, to Arabic \emph{iʿtimād}, \textquotedblleft{}reliance, leaning on\textquotedblright{}.
\end{cousinweb}

\begin{grammarlens}[title={What you've built}]
One more word for saying what happens if.
\end{grammarlens}

\subsection*{Your turn}
\begin{itemize}
\raggedright
  \item \emph{Say it:} \emph{etimād}
  \item \emph{Say it:} \emph{etimād}, once more
  \item \emph{Say it:} \emph{agar etimād hai to āp zarūr jīteṅ ge.} — If you have confidence, you will certainly win.
  \item \emph{Recall:} say the Urdu for a symptom, then the Urdu for a nagging fear, then say \emph{etimād} again
\end{itemize}

\begin{tcolorbox}[breakable,colback=teal!4,colframe=teal!35!black,title={Before you move on}]
What does \ur{اعتماد} mean? (\textquotedblleft{}Confidence; trust\textquotedblright{}.) Say it once more.
\end{tcolorbox}
